\section{Comparing observations with the same innovations}
\label{sec:coupling}

The innovations constructed in the preceding section are independent
Brownian motions, but that fact alone does not show that they determine
the observations. We now turn this question into a comparison of two
observation processes. We sample them independently conditional on the
same \emph{entire} innovation array, prove that they can still be driven
by the same Brownian motions in their joint filtration, and derive an
equation for their discrepancy. This is the conditional-copy
compatibility method developed by Kurtz~\cite{Kurtz2014}; we prove the
specific compatibility and transport statements needed here.

Throughout this section, $d\ge3$, $b=d-1$, and
$0<\theta=\tanh\beta\le b^{-1/2}$. We use the continuous causal
observation functions selected in Proposition~\ref{prop:innovations}.
With the selected means, the innovation map
\[
 W_v(t)=X_v(t)-\int_0^t m_v(s)\dd s
\]
is a Borel function of $X$ whose restriction to $[0,t]$ uses only
observations up to $t$. By Section~\ref{sec:innovations} and a countable
intersection over vertices, this selected-mean innovation map agrees
with the original raw innovation map as a whole-path map almost surely
under the observation law. We use this selected version throughout.
The field equation~\eqref{eq:field-sde} and
the integrated row bound~\eqref{eq:row-energy} have already been
proved in the observation filtration. No inverse map from $W$ to $X$
is assumed in what follows.
For drift integrals involving $A_v$, we use the finite-or-zero drift
convention fixed in Section~\ref{sec:innovations}.

\subsection{Conditioning on the complete innovations}

Write
\[
 \begin{gathered}
 C_0([0,\infty),\mathbb R)
   =\{x\in C([0,\infty),\mathbb R):x(0)=0\},\\
 \mathsf X=\mathsf W=C_0([0,\infty),\mathbb R)^{\mathcal V}.
 \end{gathered}
\]
Give the path space its locally uniform Polish topology and the arrays
the countable product topology. The subscript $0$ specifies the value
at time zero and imposes no condition at infinity.
Let $\mathsf Q$ be the joint law of $(X,W)$ from the
planted observation construction. Its $W$-marginal is product Wiener
measure, denoted by $\mathsf w$. These are standard Borel spaces, so
there is a regular conditional probability $\mathsf K(w,\dd x)$ for
$X$ given $W=w$. Define a probability law on the three path arrays by
\begin{equation}\label{eq:copy-law}
 \widehat{\mathsf Q}(\dd w,\dd x,\dd x')
   =\mathsf w(\dd w)\,\mathsf K(w,\dd x)\,\mathsf K(w,\dd x').
\end{equation}
Write $X^1,X^2,W$ for its coordinates. Each pair $(X^j,W)$ has law
$\mathsf Q$, while $X^1$ and $X^2$ are independent conditional on
$W$. The conditioning in~\eqref{eq:copy-law} uses every coordinate
of $W$ at every time, not only the portion already observed.

For $t\ge0$, let $\mathcal H_t^0$ be the raw sigma field generated by
$X^1_v(s),X^2_v(s),W_v(s)$ for $v\in\mathcal V$ and $0\le s\le t$,
and let $(\mathcal H_t)$ be its usual augmentation. All posterior
fields and coefficients attached to copy $j$ below are obtained by
applying the same causal observation functions to $X^j$.

\begin{proposition}[Compatible conditional copies]
\label{prop:compatible-copies}
The law~\eqref{eq:copy-law} is invariant under the simultaneous
action of each fixed automorphism of $T_d$. Every finite subvector of
$W$ is standard vector Brownian motion in $(\mathcal H_t)$. For each
copy $j$, both equations~\eqref{eq:magnetization-sde} and
\eqref{eq:field-sde} hold in this filtration, with the coefficients
of that copy and the common Brownian motions $W$. Their infinite
Brownian sums converge in the maximal $L^2$ norm on every finite
time interval.
\end{proposition}

\begin{proof}
We first prove that the future of $W$ is not exposed by either
observation past. Fix a deterministic $t\ge0$, write $W_{\le t}$ for
the restriction of $W$ to $[0,t]$, and define its future increment
array by
\[
 Z^t_v(s)=W_v(t+s)-W_v(t),\qquad s\ge0.
\]
In the original law $\mathsf Q$, $Z^t$ is independent of the raw
observation past $\mathcal F_t^0$. Indeed, the Brownian property in
Proposition~\ref{prop:innovations} gives the assertion for every
finite set of sites and future times. Rational times generate the
Borel sigma field of a continuous path; a monotone-class argument
over the countably many sites and rational times gives the assertion
for the whole increment array. Causality also makes $W_{\le t}$
$\mathcal F_t^0$-measurable.

It follows that, for every bounded Borel function $F$ of the
observation past,
\begin{equation}\label{eq:past-compatibility}
 \E_{\mathsf Q}[F(X_{\le t})\mid W]
   =\E_{\mathsf Q}[F(X_{\le t})\mid W_{\le t}]
   \quad\text{almost surely}.
\end{equation}
To verify the identity, denote the right side by $J(W_{\le t})$ and
test against $u(W_{\le t})v(Z^t)$, with bounded Borel $u,v$.
Independence of $Z^t$ from $\mathcal F_t^0$ factors $\E v(Z^t)$ out
of the expectation containing $F u$. It does the same for $J u$,
and the remaining expectations agree by the definition of $J$.
These products generate $\sigma(W)$: its past together with $Z^t$
reconstructs every path value of $W$. This proves
\eqref{eq:past-compatibility}.

Conditional independence in~\eqref{eq:copy-law} gives, for bounded
past functions $F,G$,
\[
 \E_{\widehat{\mathsf Q}}
   [F(X^1_{\le t})G(X^2_{\le t})\mid W]
 =\E_{\mathsf Q}[F(X_{\le t})\mid W]
  \E_{\mathsf Q}[G(X_{\le t})\mid W].
\]
Both factors on the right depend only on $W_{\le t}$ by
\eqref{eq:past-compatibility}. Product Wiener measure makes $Z^t$
independent of $W_{\le t}$, and hence
\begin{align*}
 &\E_{\widehat{\mathsf Q}}
   [u(W_{\le t})F(X^1_{\le t})G(X^2_{\le t})v(Z^t)]\\
 &\qquad=\E_{\widehat{\mathsf Q}}
   [u(W_{\le t})F(X^1_{\le t})G(X^2_{\le t})]
       \E_{\mathsf w}[v(Z^t)].
\end{align*}
Another monotone-class argument shows that $Z^t$ is independent of
the full sigma field $\mathcal H_t^0$. Thus every finite subvector
of $W$ is Brownian in the raw joined filtration.

Each $W_v$ is a continuous square-integrable martingale in the raw
joined filtration and is $L^2$-continuous because its marginal law is
Wiener measure. Lemma~\ref{lem:continuous-augmentation} therefore makes
it a martingale in $(\mathcal H_t)$. The quadratic covariations,
determined by the Brownian paths, remain
$[W_u,W_v]_t=\mathbf 1_{\{u=v\}}t$. L\'evy's martingale
characterization applied to each finite subvector proves the claimed
Brownian property in $(\mathcal H_t)$.

We next check invariance. The selected field functions commute with
graph relabeling, so the innovation map does too. Since the observation
law is invariant, $\mathsf Q$ is invariant under
$(X,W)\mapsto(gX,gW)$. For one fixed automorphism $g$, uniqueness of
regular conditional probabilities now gives
\begin{equation}\label{eq:kernel-equivariance}
 \mathsf K(gw,\cdot)=g_*\mathsf K(w,\cdot)
       \quad\text{for $\mathsf w$-almost every }w.
\end{equation}
To obtain equality as probability measures, first check it for a
countable determining algebra generating the Borel sets of $\mathsf X$
and intersect the resulting full-measure sets. Applying this identity
to both factors in~\eqref{eq:copy-law} proves invariance of the
joined law. This argument is made separately for each fixed $g$.

It remains to justify the stochastic equations in the enlarged
filtration. We give the approximation argument because an It\^o
integral cannot simply be identified across filtrations from a
formal equality of pair laws. For either coefficient
$k_{vu}=c_{vu}$ or $k_{vu}=C_{vu}$, fix $v,u$ and a finite horizon
$T$. Lemma~\ref{lem:observation-step-approximation}, applied to the
$X$-marginal of $\mathsf Q$, supplies bounded causal Borel elementary
processes
\[
 G_n(t,X)=\sum_i \gamma_{n,i}(X_{\le t_{n,i}})
       \mathbf 1_{(t_{n,i},t_{n,i+1}]}(t),
 \qquad
 \int_0^T\E_{\mathsf Q}|G_n-k_{vu}|^2\dd t\longrightarrow0,
\]
where each $\gamma_{n,i}$ is a bounded Borel function of finitely many
observations at time $t_{n,i}$, viewed as a function of the indicated
past.

The elementary integral against $W_u$ is the Borel path functional
\[
 I_n(t;X,W)=\sum_i\gamma_{n,i}(X_{\le t_{n,i}})
    \bigl(W_u(t\wedge t_{n,i+1})-W_u(t\wedge t_{n,i})\bigr).
\]
The Brownian isometry and Doob's maximal inequality show that these
integrals converge in $L^2(C([0,T]))$ to the original It\^o integral.
A deterministic subsequence converges uniformly almost surely; its
limit on the Borel set of uniformly Cauchy inputs, with the zero path
as default, is a Borel version of that integral as a function of
$(X,W)$. The same subsequence has that pathwise limit in each joined
pair, since $(X^j,W)$ has law $\mathsf Q$. On the other hand the same
isometry and maximal inequality in $(\mathcal H_t)$ show that
$I_n(\cdot;X^j,W)$ converges to the joined It\^o integral of
$k_{vu}(\cdot,X^j)$. The two limits agree. This proves transport of
each individual integral, including its continuous path version.

For either coefficient family the integrated squared row is finite:
this is~\eqref{eq:row-energy} for $c$, and follows for $C$ also from
$C_{vu}=(1-m_v^2)c_{vu}$. If $F\subset F'$ are finite vertex sets,
the joined integral tails satisfy
\[
 \E_{\widehat{\mathsf Q}}\sup_{t\le T}
  \left|\sum_{u\in F'\setminus F}\int_0^t
                k_{vu}(s,X^j)\dd W_u(s)\right|^2
 \le4\sum_{u\in F'\setminus F}\E_{\mathsf Q}
                    \int_0^T k_{vu}(s,X)^2\dd s.
\]
The right side tends to zero along an exhaustion. The drift at any
fixed time is a Borel observation functional, obtained by time
integration; its absolute integral is finite almost surely because
$|m_v|\le1$ and $\E\int_0^T A_v\dd t<\infty$. Assigning zero to
exceptional nonfinite integral values supplies a Borel scalar version
if needed. Transport the original equations with finite partial
Brownian sums at rational times and let the partial sums increase:
their residuals converge in probability under $\mathsf Q$, hence
under each identical joined pair law, while their joined integral
limits have just been identified. The equations therefore hold at
all rational times in the joining. Continuity of the fields and
Brownian sums, and absolute continuity of the drift integrals, extend
them to all times. Taking countable intersections over vertices and
integer horizons finishes the transport assertion.
\end{proof}

The purpose of this construction can already be stated precisely.
For each vertex $v$ and nonnegative rational $q$,
$\E_{\mathsf Q}X_v(q)^2=q+q^2<\infty$. Conditional independence gives
\begin{equation}\label{eq:copy-coordinate-variance}
 \E_{\widehat{\mathsf Q}}
       [(X^1_v(q)-X^2_v(q))^2\mid W]
   =2\Var_{\mathsf Q}(X_v(q)\mid W)
       \quad\text{almost surely}.
\end{equation}
Thus equality of the two observation paths would make every rational
coordinate have zero conditional variance. These countably many
coordinates determine a continuous path configuration. The final
section will turn that implication into a Borel recovery map after
the comparison has proved equality; no equality has been assumed here.

\subsection{The field discrepancy and its variance}

From now on unprimed quantities refer to the first copy, primed
quantities to the second, and $\Delta$ means their difference. All
expectations in the comparison are under $\widehat{\mathsf Q}$.
Every one-copy estimate proved earlier remains valid because the
marginals have not changed.

Fix $T>0$. In each pair the defining innovation identity holds as an
identity of continuous paths. Subtracting the two identities gives
\[
 \Delta X_v(t)=\int_0^t\Delta m_v(s)\dd s,
 \qquad |\Delta X_v(t)|\le2t.
\]
Together with $|h_v-X_v|,|h'_v-X'_v|\le d\beta$, this gives
\begin{equation}\label{eq:bounded-discrepancy}
 |\Delta h_v(t)|\le H_T:=2T+2d\beta
       \qquad(v\in\mathcal V, 0\le t\le T)
\end{equation}
on one full-probability event. The field bounds and identities hold
simultaneously because the chosen fields are continuous and the
vertices are countable.

Choose a vertex $o$ only for taking expectations, and put
\[
 V(t)=\E[(\Delta h_o(t))^2].
\]
Invariance of the joined law makes the value independent of the
choice of $o$. It is nonnegative, bounded by $H_T^2$ on $[0,T]$,
and $V(0)=0$. We need the full variance equation, including its drift.

\begin{lemma}[Variance of the discrepancy]
\label{lem:copy-variance}
The function $V$ is absolutely continuous on every finite interval.
For almost every $t$,
\begin{equation}\label{eq:variance}
 V'(t)=\sum_{u\in\mathcal V}\E[(\Delta c_{ou}(t))^2]
       +2\E[\Delta h_o(t)(m_o(t)A_o(t)-m'_o(t)A'_o(t))].
\end{equation}
\end{lemma}

\begin{proof}
Subtract the two field equations in the common filtration. The
martingale part is $M_o=\sum_u\int\Delta c_{ou}\dd W_u$, and, outside
the $\dd t\,\dd\widehat{\mathsf Q}$-null set where a row diverges,
\[
 \sum_u(\Delta c_{ou})^2\le2(A_o+A'_o).
\]
The integrated expectation of the right side is finite. Thus the
sum defining $M_o$ converges in maximal $L^2$, and its bracket is
\[
 \langle M_o\rangle_t
       =\int_0^t\sum_u(\Delta c_{ou}(s))^2\dd s.
\]
This follows first for finite sums from the zero cross-covariations
of distinct Brownian coordinates, then for the limit by the
$L^1$ convergence of the integrated squared row tails. It\^o's
formula for $(\Delta h_o)^2$ gives
\begin{align*}
 (\Delta h_o(t))^2={}&2\int_0^t\Delta h_o(s)\dd M_o(s)
  +\int_0^t\sum_u(\Delta c_{ou}(s))^2\dd s\\
 &+2\int_0^t\Delta h_o(s)
          (m_o(s)A_o(s)-m'_o(s)A'_o(s))\dd s.
\end{align*}
The first integral is a square-integrable martingale: by
\eqref{eq:bounded-discrepancy}, its expected bracket is at most
$8H_T^2\E\int_0^T(A_o+A'_o)\dd s$. The absolute values of the
other two integrands have expectations bounded respectively by
$2\E(A_o+A'_o)$ and $2H_T\E(A_o+A'_o)$, which are integrable in
time. Taking expectations therefore gives an integral of an
$L^1([0,T])$ function. This proves absolute continuity and its
almost-everywhere derivative~\eqref{eq:variance}.
\end{proof}

The coefficients in this equation are products along directed paths.
We first control the difference of one factor. For a directed edge
$e=(i,j)$ define
\begin{equation}\label{eq:log-discrepancy-def}
 D_e(t)=\log a(h_{j\to i}(t))-
        \log a(h'_{j\to i}(t)),
 \qquad S(t)=\E[D_e(t)^2].
\end{equation}
Transitivity on directed edges makes $S$ independent of $e$.
The two cavity identities at the ends of $ij$ give
\[
 |\Delta h_{j\to i}|
     \le|\Delta h_j|+\theta|\Delta h_{i\to j}|,
 \qquad
 |\Delta h_{i\to j}|
     \le|\Delta h_i|+\theta|\Delta h_{j\to i}|.
\]
Substituting the second inequality into the first yields
\begin{equation}\label{eq:component-difference}
 |\Delta h_{j\to i}|
  \le\frac{|\Delta h_j|+\theta|\Delta h_i|}{1-\theta^2}.
\end{equation}
In particular, on $[0,T]$ it is bounded by $H_T/(1-\theta)$;
Minkowski's inequality and invariance also give
$\|\Delta h_{j\to i}\|_2\le\sqrt V/(1-\theta)$.

Direct differentiation gives
\begin{equation}\label{eq:log-derivative}
 (\log a)'(z)
   =-\frac{2(1-\theta^2)\tanh z}{1-\theta^2\tanh^2z},
 \qquad |(\log a)'(z)|\le2|\tanh z|\le2.
\end{equation}
Here the first inequality uses
$1-\theta^2\tanh^2z\ge1-\theta^2$. We conclude that
\begin{equation}\label{eq:basic-log-loss}
 |D_e(t)|\le\frac{2H_T}{1-\theta},
 \qquad S(t)\le\frac{4}{(1-\theta)^2}V(t)
       \quad(0\le t\le T).
\end{equation}
Also $(\log a)(0)=0$, so
$|\log a(h_{j\to i}(t))|\le2|h_{j\to i}(t)|
\le2(|X_j(t)|+b\beta)$. Since $\E X_j(t)^2=t+t^2$, each one-copy
log factor is square integrable. Equality of the marginal laws then
gives $\E D_e(t)=0$. These facts will permit the later spatial
cancellation estimate to be applied to this particular edge observable.

\subsection{Reducing the comparison to path estimates}

For a directed simple path $v_0,\ldots,v_r$, $r\ge1$, let
$e_i=(v_{i-1},v_i)$. At a common time $t$, set
\begin{equation}\label{eq:weighted-path-definitions}
 \begin{aligned}
 P_r&=\prod_{i=1}^r a(h_{v_i\to v_{i-1}}),
 &Z_r&=\sum_{i=1}^rD_{e_i},\\
 Q_r&=\E[(P_r^2+(P'_r)^2)Z_r^2],
 &L_r&=\E[(\Delta h_{v_0})^2P_r^2].
 \end{aligned}
\end{equation}
The orientation matters: every factor uses the component at the next
vertex excluding the preceding one. Invariance makes these expectations
depend only on $r$ and $t$, so we suppress the particular path. Both
$P_r$ and $P'_r$ lie in $(0,1]$, and
$\log P_r-\log P'_r=Z_r$. The mean-value theorem, first for
$e^x$ and then for $e^{2x}$, gives
\begin{equation}\label{eq:product-difference}
 \E[(P_r-P'_r)^2]\le2Q_r,
 \qquad
 |P_r^2-(P'_r)^2|
       \le2(P_r^2+(P'_r)^2)|Z_r|.
\end{equation}
For example, the derivative of $e^x$ on the segment joining
$\log P_r$ and $\log P'_r$ is at most $P_r+P'_r$; squaring and
using $(P_r+P'_r)^2\le2(P_r^2+(P'_r)^2)$ proves the first bound.
The derivative of $e^{2x}$ on the same segment is at most
$2(P_r^2+(P'_r)^2)$, proving the second.

There are $n_r=db^{r-1}$ vertices at distance $r\ge1$ from $o$.
The diagonal coefficients are $c_{oo}=c'_{oo}=1$, and the
coefficient along any positive-length path is $\theta^rP_r$.
Tonelli's theorem and~\eqref{eq:product-difference} therefore bound
the noise term in~\eqref{eq:variance} by
\begin{equation}\label{eq:noise-bound}
 \sum_u\E[(\Delta c_{ou})^2]
       \le2\sum_{r\ge1}n_r\theta^{2r}Q_r.
\end{equation}

The drift contributes a second, different path sum:
\begin{equation}\label{eq:drift-bound}
 \left|\E[\Delta h_o(m_oA_o-m'_oA'_o)]\right|
 \le V+\sum_{r\ge1}n_r\theta^{2r}
          \left(L_r+2\sqrt{2V}\sqrt{Q_r}\right).
\end{equation}
To prove this, write
$m_oA_o-m'_oA'_o=(m_o-m'_o)A_o+m'_o(A_o-A'_o)$.
Since $\tanh$ is $1$-Lipschitz, the first part is bounded by
\[
 \E[(\Delta h_o)^2A_o]
       =V+\sum_{r\ge1}n_r\theta^{2r}L_r.
\]
For the second, use $|m'_o|\le1$. The diagonal terms of $A_o-A'_o$
cancel. For each remaining path, put $w_r=P_r^2+(P'_r)^2\le2$.
The second inequality in~\eqref{eq:product-difference} and weighted
Cauchy--Schwarz give
\begin{align*}
 \E[|\Delta h_o|\,|P_r^2-(P'_r)^2|]
 &\le2\sqrt{\E[(\Delta h_o)^2w_r]}
          \sqrt{\E[w_rZ_r^2]}\\
 &\le2\sqrt{2V}\sqrt{Q_r}.
\end{align*}
Summing proves~\eqref{eq:drift-bound}. These inequalities may first
be read for finite spatial sums and then with nonnegative extended
right sides; the estimates at their subsequent uses make those
right sides finite. In particular no independence between a path
product and its discrepancy is used. The square-root sum in the
drift must be controlled along with the noise sum.

\begin{proposition}[Uniqueness below the threshold]
\label{prop:strict-uniqueness}
If $b\theta^2<1$, then $V(t)=0$ for every $t\ge0$.
\end{proposition}

\begin{proof}
Since the products in~\eqref{eq:weighted-path-definitions} are at
most one, $L_r\le V$. Cauchy--Schwarz in the sum $Z_r$ gives
\[
 Q_r\le2\E Z_r^2
       \le2r\sum_{i=1}^r\E D_{e_i}^2=2r^2S.
\]
Use this with~\eqref{eq:basic-log-loss} in the noise and drift
bounds. The contributions are bounded by $V$ times sums of
$r^2$, $1$, and $r$, respectively, against
\[
 n_r\theta^{2r}=\frac db(b\theta^2)^r.
\]
All three series converge when $b\theta^2<1$. Equation
\eqref{eq:variance} therefore gives $V'\le C V$ almost everywhere
on $[0,T]$, for a finite constant depending on the fixed parameters.
Absolute continuity and $V(0)=0$ imply $V=0$ there by Gronwall's
inequality. The horizon $T$ was arbitrary.
\end{proof}

At equality, $n_r\theta^{2r}=d/b$, so the geometric decay in this
proof disappears. The next two sections will replace the elementary
$r^2S$ loss by spatial cancellation of order $rS$ and attenuation
on the length scale $\omega(t)^{-1}$. Their leading noise estimate
has the form $Q_r\lesssim rS(t)e^{-c\omega(t)r}$. Near time zero,
where $\omega(t)=\sqrt t$, its sum has size
\[
 S(t)\sum_{r\ge1}r e^{-c\sqrt t\,r}
   =S(t)\frac{e^{-c\sqrt t}}{(1-e^{-c\sqrt t})^2}
   \asymp\frac{S(t)}{t}.
\]
Thus even those two estimates, combined only with $S\le CV$, would
leave a nonintegrable coefficient at time zero. The final comparison
will improve the relation between $S$ and $V$ using small component
fields, while also retaining the square-root drift sum. This is the
remaining obstruction at the critical parameter.
